\section{Exact boundaries and counterexamples}
\label{sec:boundaries}

The hypotheses identify useful mathematical structure. Several of them
cannot be discarded merely because each finite computation terminates.
The following elementary examples make the boundaries testable.

\subsection{Separately contracting implementations can disagree}

On \([0,1]\cap\Q\), let \(E_p(x)=0\) for even \(p\) and
\(E_p(x)=1\) for odd \(p\). Each map is a constant, hence a
contraction of factor zero with zero fixed-point residual. Their
fixed points alternate. No regular name follows.

The mixed-precision estimate (SG3) fails: for large adjacent
precisions its right side at equal input is \(2^{-p}+2^{-p-1}\),
while its left side is one. Thus cross-precision coherence is a
separate input property, and cannot be inferred from contraction of
each implementation in isolation.

\subsection{Vanishing commutation defects do not overcome a closing gap}

Take \(D_n=[0,1]\cap\Q\), identity restrictions,
\(\eps_n=2^{-n-2}\), and
\[
 F_n(x)=(1-\eps_n)x+\eps_nc_n,
 \qquad
 c_n=\begin{cases}0&n\text{ even},\\1&n\text{ odd}.
 \end{cases}
\]
Each map has the exact rational fixed point \(y_n=c_n\), and
contraction factor \(q_n=1-\eps_n<1\). Direct subtraction gives
\[
 \sup_{0\le x\le1}|F_n(x)-F_{n+1}(x)|\le\eps_n.
\]
For example, when \(c_n=0,c_{n+1}=1\), the difference is
\(-\eps_nx/2-\eps_n/2\); the other parity is analogous.
The commutation defect tends to zero with an explicit dyadic rate,
but the exact fixed points alternate between zero and one.

There is no contradiction with (R2): the defect divided by the
contraction gap is \(\eps_n/(1-q_n)=1\). It does not have a
vanishing tail. This example shows why any theorem that omits the gap
factor would be false even for one-dimensional affine maps.

\subsection{A local numerical certificate is not a circular proof rule}

A system of equations may have cycles in its dependency graph. The
present method handles a subclass by replacing an implicit solution
with a finite sequence of explicit updates and proving an error
invariant for that sequence. The proof of each update estimate is
already available before the iteration theorem is applied.

This does not license a proof graph in which a claim assumes itself
through another unproved claim. No equation such as ``certificate
\(c_1\) is valid because \(c_2\) is valid, and conversely'' is
accepted merely because its dataflow diagram is finite. An unrelated
cyclic certificate system would need a sound induction, descent or
guardedness argument of its own. Thus the open cyclic-execution item
in the repository is narrowed for numerical feedback with small-gain
data, not solved in its full generality.

\subsection{No total primitive recursive universal evaluator}

Suppose a total PR function \(U(e,n)\) evaluated every unary PR
program at its code \(e\). Then \(g(n)=U(n,n)+1\) would itself be
PR, with some code \(a\). The defining property at \((a,a)\)
would give
\(U(a,a)=g(a)=U(a,a)+1\), a contradiction. Consequently
``all producer codes are PR'' cannot be implemented by adding a
single total PR evaluator for them.

Our two legitimate uses of uniformity are different. A fixed PR
body can have arbitrarily many numerical parameters. Alternatively,
a fixed bounded syntax such as finite rational polynomials has its
own PR evaluator. General program traces can also be checked with a
supplied fuel bound, but the claim that the bound is sufficient is
then an explicit obligation.

\subsection{Metastability and convergence remain different outputs}

A finite stable window does not by itself supply a tail modulus or
a coherent name. The proof-mining literature explicitly distinguishes
rates of convergence, finite learning information, and rates of
metastability \cite{kohlenbachsafarik}. The current construction obtains
a full modulus because the gap and the error schedule give a finite
geometric estimate. A theorem lacking those data may still have
valuable metastable consequences; it must retain that weaker output
instead of applying the regular-name theorem without its hypotheses.

\subsection{Cardinality is not a complete enumeration}

Even a bound of one possible witness does not generally produce a
search cutoff. For a program code \(e\), the relation ``\(t\) is
the first halting time of program \(e\)'' is decidable by bounded
simulation and has at most one witness. A uniform computable bound
containing that witness whenever it exists would decide halting by
simulation up to the returned bound. The usual diagonal argument
rules this out.

This elementary example explains the logical distinction; it is not
a reduction proving anything about the particular curves in the
Faltings or Siegel packets. Their requested height bounds, complete
lists, cardinality statements and finite-tuple obstructions are
different arithmetical targets. A new proof of one target would
need its own exact hypotheses and argument. No inference of
independence, unprovability in iPRA, or impossibility of every
constructive approach is made here.

\subsection{Failure of a certificate is inconclusive}

A failed weighted-row test does not prove nonexistence of a solution.
Weights, a smaller box, a different preconditioner, or another
mathematical method may succeed. Likewise, a passing local test does
not establish global continuation, uniqueness outside its region,
or behavior at a singular parameter. The numerical inequalities
state the scope of the result exactly.
